% TOP_PROBLEM: {"id": 30000030, "title": "Cannon's Conjecture", "category_id": 6, "category": {"id": 6, "name": "geometry", "display_name": "Geometry", "description": "Euclidean and non-Euclidean geometry, geometric structures.", "slug": "geometry", "order_index": 6, "created_at": "2026-07-31T15:26:25.671Z"}, "status": "open"}
% FIRST_POSTED: null
% !TeX program = lualatex
% Compile twice: lualatex 30000030.tex
% All references and prose are embedded; no BibTeX database or external inputs.
\documentclass[11pt,a4paper]{article}
\usepackage[margin=24mm,headheight=14pt]{geometry}
\usepackage{fontspec}
\setmainfont{Latin Modern Roman}
\setsansfont{Latin Modern Sans}
\setmonofont{Latin Modern Mono}
\usepackage{amsmath,amssymb,mathtools}
\usepackage{unicode-math}
\setmathfont{Latin Modern Math}
\usepackage{microtype}
\usepackage{enumitem,xurl,needspace}
\usepackage[dvipsnames]{xcolor}
\usepackage{hyperref}
\hypersetup{unicode=true,colorlinks=true,linkcolor=MidnightBlue,urlcolor=MidnightBlue,
 pdftitle={Research attempt: Problem 30000030},pdfauthor={Research notebook},bookmarksnumbered=true}
\usepackage{bookmark}
\usepackage{tocloft}
\setlength{\cftsecnumwidth}{3em}
\cftsetpnumwidth{2.5em}
\cftsetrmarg{4em}
\usepackage{titlesec}
\titleformat{\section}{\Large\bfseries\raggedright}{\thesection}{0.7em}{}
\usepackage{fancyhdr}
\pagestyle{fancy}\fancyhf{}
\fancyhead[L]{\small\sffamily Open Problems Math | Research notebook}
\fancyhead[R]{\small Problem 30000030 | 27 September 2026}
\fancyfoot[C]{\thepage}
\setlength{\parindent}{0pt}
\setlength{\parskip}{3pt plus 1pt}
\setlength{\emergencystretch}{3em}
\setcounter{tocdepth}{1}
\setcounter{secnumdepth}{1}
\setlist[itemize]{leftmargin=3em,itemsep=5pt,topsep=4pt,parsep=0pt}
\allowdisplaybreaks
\newcommand{\N}{\mathbb{N}}
\newcommand{\Z}{\mathbb{Z}}
\newcommand{\Q}{\mathbb{Q}}
\newcommand{\R}{\mathbb{R}}
\newcommand{\C}{\mathbb{C}}
\newcommand{\F}{\mathbb{F}}
\newcommand{\A}{\mathbb{A}}
\newcommand{\PP}{\mathbb P}
\newcommand{\E}{\mathbb{E}}
\newcommand{\eps}{\varepsilon}
\newcommand{\dd}{\,\mathrm d}
\newcommand{\sref}[2]{\hyperlink{source:#1:#2}{[S#2]}}
\newcommand{\eref}[2]{\hyperlink{extra:#1:#2}{[E#2]}}
% Turn the next switch off for a shorter reading copy. The quotations preserve
% every exactTarget field, including qualifications omitted from a short summary.
\newif\ifshowcatalogue
\showcataloguetrue
\newcommand{\cataloguescope}[1]{\ifshowcatalogue
 \paragraph{Exact scope in the supplied catalogue.}
 \begin{quote}\small #1\end{quote}\fi}
\usepackage{etoolbox}
\AtBeginEnvironment{itemize}{\small}
\numberwithin{equation}{section}
% Standalone export. Original research content preserved.
\begin{document}
\setcounter{section}{0}
% ================================================================
% problemId: 30000030
% Canonical problem ID: 30000030
% exactTarget SHA-256: 28a0b3ddf25701e2015369688334ad40f8a63b668a17b30d1662e3f8ebb63d97
\clearpage
\section*{\texorpdfstring{Cannon's Conjecture}{Cannon's Conjecture}}\label{30000030}
\subsection{Definitions and mathematical statement}
A finitely generated group is Gromov-hyperbolic if a Cayley graph has uniformly thin geodesic triangles. Its boundary consists of asymptotic classes of geodesic rays, with the standard visual topology. The conjecture is
\[
 \partial G\cong S^2\quad\Longrightarrow\quad
 G\curvearrowright\mathbb H^3\text{ properly, cocompactly, and isometrically}.
\]
Properness allows finite stabilizers in the presence of torsion. Equivalently the group is virtually a cocompact Kleinian group in the source's convention. The hypothesis is about the boundary's topology, not merely its dimension.
\par\smallskip\textit{Formulation and concept references:} \sref{30000030}{1}.
\cataloguescope{Prove or disprove that every Gromov-hyperbolic group whose boundary at infinity is homeomorphic to the two-sphere acts properly discontinuously, cocompactly, and isometrically on hyperbolic three-space, equivalently is virtually a cocompact Kleinian group.}
\subsection{Short English statement}
Must a hyperbolic group with a two-sphere boundary come from a cocompact hyperbolic three-dimensional geometry?
% BEGIN RESEARCH ADDITION 140
\subsection{Research attempt}

\paragraph{Current frontier.}
The analytic route asks whether a visual boundary metric is quasisymmetric
to the round sphere. Bonk--Kleiner explain the relation to Cannon's geometric
action conclusion and prove uniformization for Ahlfors $2$-regular, linearly
locally contractible metric spheres. Topological dimension two is not
Ahlfors $2$-regularity. Their graph-modulus criteria offer another route
without assuming that regularity at the outset. \eref{30000030}{1}
Markovi\'c's alternative criterion uses enough quasiconvex surface subgroups
to separate boundary point pairs; their existence is an extra requirement,
not a consequence proved by the criterion. \eref{30000030}{2}
The geometric-measure reformulation of Cheetham-West--Nolte is likewise
not used here as an automatic regularity theorem. \eref{30000030}{3}

\paragraph{Attack route.}
Transfer quadratic energies from controlled planar model graphs to finer
boundary approximations. The aim is to replace ``approximately round'' by
verifiable path lengths and congestion bounds. Hyperbolicity supplies
multiscale structure, but it does not by itself provide the scale-uniform
comparison established conditionally below. The alternative surface-subgroup
route would require an actual construction of separating subgroups.

\paragraph{Attempt: a finite network-transfer lemma.}
For a finite connected undirected graph $\mathcal G$ and nonempty disjoint
terminal sets $A,B$, define
\[
 \operatorname{cap}_{\mathcal G}(A,B)
 =\min_{f|_A=0,\ f|_B=1}\sum_{\{v,w\}\in E(\mathcal G)}(f(v)-f(w))^2.
\]
Clipping to $[0,1]$ cannot increase energy; a minimum thus exists on a compact
feasible set. Suppose $\varphi:V(\mathcal H)\to V(\mathcal G)$ takes
terminals to corresponding terminals and represents every edge
$e=\{v,w\}$ of $\mathcal H$ by a path $P_e$ joining $\varphi(v),\varphi(w)$.
If paths have length at most $L$ and each edge of $\mathcal G$ belongs to at
most $M$ paths, counting multiplicity, then
\begin{equation}\label{30000030:transfer}
 \operatorname{cap}_{\mathcal H}\le LM\operatorname{cap}_{\mathcal G}.
\end{equation}
For any admissible $f$ on $\mathcal G$, Cauchy--Schwarz gives
\[
 |f(\varphi(v))-f(\varphi(w))|^2
 \le |P_e|\sum_{a\in P_e}|df(a)|^2.
\]
Sum over $e$, use congestion, then minimize. A reverse map with parameters
$L',M'$ gives
\[
 (LM)^{-1}\operatorname{cap}_{\mathcal H}
 \le\operatorname{cap}_{\mathcal G}
 \le L'M'\operatorname{cap}_{\mathcal H}.
\]
More precisely $LM$ can be replaced by
$\max_a\sum_{e:a\in P_e}|P_e|$, with multiplicities retained. This weighted
congestion may be smaller when few paths are long. No limiting argument is
involved in this lemma.

\paragraph{An exactly soluble model.}
Let $\mathcal R_{m,n}$ have vertices $(i,j)$, $0\le i\le m$, $1\le j\le n$,
nearest-neighbor edges, and terminals $i=0,m$. Each of its $n$ horizontal
rows satisfies
\[
 \sum_{i=0}^{m-1}(f(i+1,j)-f(i,j))^2\ge\frac1m.
\]
Hence capacity is at least $n/m$. The potential $f(i,j)=i/m$ attains this
bound and has zero vertical energy. Thus
\begin{equation}\label{30000030:grid}
 \operatorname{cap}_{\mathcal R_{m,n}}=n/m.
\end{equation}
All graph degrees are at most four, yet capacity can tend to zero or
infinity. Planarity and bounded valence do not supply uniform conformal
control.

The issue has a metric realization. On a coordinate square, let
\[
 d_\alpha((x,y),(x',y'))=|x-x'|+|y-y'|^\alpha,
 \qquad 0<\alpha<1.
\]
This has the usual topology; balls contract under coordinate contraction.
A comparable-mesh grid at scale $\epsilon$ has horizontal spacing
$\asymp\epsilon$ and vertical spacing $\asymp\epsilon^{1/\alpha}$.
Thus $m\asymp\epsilon^{-1}$ and $n\asymp\epsilon^{-1/\alpha}$, giving
\[
 \operatorname{cap}_{\rm left,right}\asymp
 \epsilon^{1-1/\alpha}\longrightarrow\infty,
\]
while transverse capacity tends to zero. This illustrates the anisotropic
local obstruction discussed in Bonk--Kleiner, not a constructed boundary of
a hyperbolic group. \eref{30000030}{1}

\paragraph{Attempt to close the multiscale step.}
If boundary approximation graphs could be compared in both directions to
round-sphere approximations as above, uniformly in scale and with controlled
terminal continua and relative distances, they would inherit round capacity
bounds. To invoke the actual uniformization criteria one must additionally
compare this edge energy to the source's graph modulus and verify its
approximation hypotheses. Bounded valence allows local energy comparisons;
it does not supply terminal geometry or compatible boundary maps. The
relevant criteria are in Sections 10--11 of the cited manuscript.
\eref{30000030}{1}

A naïve level-by-level construction is insufficient. Composing two path
maps multiplies available length and congestion bounds:
$L_{12}\le L_1L_2$, $M_{12}\le M_1M_2$. A cost $K>1$ per level can become
$K^j$, not an absolute constant. An energy-aware construction would need
to redistribute paths and bound weighted congestion directly across scales.
No such routing is obtained from the group axioms here. Finite checks of
many subdivisions do not establish its uniform bound at all scales.

\paragraph{Stress test.}
Good topology and bounded graph degrees coexist with the divergent capacity
example. Conversely a globally snowflaked round metric is still
quasisymmetric to the round metric, so Hausdorff dimension greater than two
alone is not a counterexample: the issue is the conformal gauge, not a raw
dimension computation. A continuous surjective limit of approximate maps
need not be a boundary homeomorphism; both directions require control.
Surface-subgroup criteria cannot be used before constructing the separating
subgroups. The exact finite-network computations in the companion script
check the lemma's normalization but supply none of these group-theoretic
inputs.

\paragraph{Outcome.}
\textbf{Outcome: Exploratory approach with unresolved gap.}
A quantitative energy-transfer lemma, exact capacity models and a multiscale
routing obstruction have been established. The missing step is a uniform,
geometrically compatible low-congestion construction for arbitrary sphere
boundaries, with every hypothesis of the uniformization criterion checked.
No new class of groups satisfying Cannon's conjecture or counterexample
group is claimed.

% END RESEARCH ADDITION 140
\subsection{Sources}
\begin{itemize}\raggedright
\item[\textnormal{[S1]}]\hypertarget{source:30000030:1}{} Expanding Thurston Maps.
\url{https://www.ams.org/bookstore/pspdf/surv-225-prev.pdf}.
{\small\textit{Catalogue formulation source.}}
% BEGIN RESEARCH REFERENCES 140
\item[\textnormal{[E1]}]\hypertarget{extra:30000030:1}{} M. Bonk and B. Kleiner,
\emph{Quasisymmetric parametrizations of two-dimensional metric spheres},
Inventiones Mathematicae 150 (2002), 127--183; manuscript
arXiv:math/0107171, Theorems 1.1--1.2, Conjecture 1.3, Sections 10--11.
\url{https://arxiv.org/pdf/math/0107171}.
\item[\textnormal{[E2]}]\hypertarget{extra:30000030:2}{} V. Markovi\'c,
\emph{Criterion for Cannon's conjecture}, Geometric and Functional Analysis
23 (2013), 1035--1061; arXiv:1205.5747.
\url{https://arxiv.org/abs/1205.5747}.
\item[\textnormal{[E3]}]\hypertarget{extra:30000030:3}{} T. Cheetham-West and A. Nolte,
\emph{Characterizing candidates for Cannon's conjecture from geometric
measure theory}, Bulletin of the London Mathematical Society 55 (2023),
1718--1725; arXiv:2207.01720v2.
\url{https://arxiv.org/html/2207.01720v2}.

% END RESEARCH REFERENCES 140
\end{itemize}

\end{document}
